% Transcription — fonds Alexandre Grothendieck, shelfmark 156-7, batch 4 (pages 61-80).
% Established from the facsimile archives/batches/156-7/batch-04.pdf (a copy of
% 156-7.pdf fetched through the project's relay; no cover sheet in the batch,
% so PDF page k = folder page 60+k), per the skill transcribe-grothendieck. The
% folder is 113 pages in six batches, transcribed in parallel; this is the
% fourth, and its neighbours (pages 41-60, 81-100) could not be relied on.
% Batches 1 and 5 of the second draft (folder 156-6, GF VI) and batch 4 of the
% first draft (156-4, GF IV: cosupport, support, « supports virtuels » Σ)
% were read for notation.
%
% Pass: Opus 5.5 (claude-opus-5-5), 2026-09-26 — first pass, unchecked against the pages by a human.
%
% Pages skipped: none. Pages summarised: none. All twenty leaves are his
% mathematics, fast drafting in ink, and all are transcribed. Page 68 is
% crossed by one long diagonal stroke and the top of page 65 (to the first
% display) and the middle of page 79 are cancelled by strokes; all are read
% and given, the cancelled blocks as \struck.
%
% His pagination 61-80 stands at the head of each leaf and coincides with the
% archivists' (noted once, on page 61). One date in his hand: « (25 juin) »,
% top left of page 69, set as \subsection{25 juin} with a \note. No other
% date in the batch.
%
% Numbered runs: the properties of supports (i)-(x) on pages 62-63 (two
% (iii) on p. 62, the first struck; a struck (v) replaced by a new (v); (x)
% read doubtfully); two conditions both numbered (iii) and a (iii') in the
% margin of page 61; conditions (i)_M-(vi)_M on page 73, (vi) with parts
% (a), (b) circled; (i)_L, (ii)_L on page 74; « Lemme » (p. 76, then called
% lemme 1 on p. 78), « Corollaire » (p. 76), « Lemme 2 » (p. 78),
% « Proposition » with Supp_P^1, Supp_P^2 and (i)-(iii) (p. 80); starred
% formulas (*) (pp. 61, 64, 75) and (**) (p. 77), (α) (p. 75), (×)
% (p. 79). Axioms cited: « At D » (p. 61, read doubtfully), Disj_Λ (p. 65,
% new), At L 1 (p. 69), At L 2 (pp. 66, 68, 73), At L 3 (p. 73), and
% « Hyp », « Hyp' » (pp. 77, 79). He cites « p. 61 » (p. 64), « p. 63, 64 »
% (p. 67) and « GF III » with « loc. cit. p. 8 » (p. 70, chapter III,
% folder 156-3), all as plain page numbers.
%
% What the batch is about. « Notion de support » in an atelier. Pages 61-64:
% Σ_F (parts of F closed for ≪ and stable under Sup for ≤) and Σ_M (parts of
% M closed for ≪) in bijection; the cosupport of A ⊂ F, cosupp A = {G | G ∥ F
% for all F ∈ A}, and supp A = cosupp cosupp A (also written Ā); the formal
% properties (i)-(x) of this Galois-type closure, F-supports Σ°_F, the
% complement ∁A = cosupp A, Inf and Sup in Σ°_F; a NB that Sup of supports
% need not be the union and distributivity may fail outside ensemblist
% ateliers. Pages 64-68: the same through M, then through any Λ ⊂ F
% satisfying Disj_Λ (disjointness detected on the Λ-shadow), in particular
% Λ = L under At L 2; the quasi-ensemblist case where ∥ on L is x ≠ y and
% Σ_L = P(L), supp A = omb(A); a generalisation to a set P ⊂ L of
% « points ». Page 69, dated 25 June: doubt that supports make sense outside
% the quasi-ensemblist case, and the example of the strata of an ordered set
% I (0-strata points, 1-strata sub-segments [a,b], ≤, ≪, compatibility (a)-(c));
% pages 70-72: finite figures there, a reference back to GF III and its
% « constructible » parts, supports attached to finite chains a_1 < … < a_n,
% and the doubt whether a set P of points exists canonically. Pages 72-80:
% a « réalisation par supports » X ↦ |X| of M in P(P) with conditions
% (i)_M-(vi)_M (reduced to lieux, (i)_L-(ii)_L); the example of a
% topological string; the question whether Σ_M ≅ P(P); Lemme (|A| ⊂ |B|
% implies supp A ⊂ supp B, with equivalence when B is a support), Corollaire
% (Σ_M → P(P) injective), Lemme 2 under « Hyp » (|cosupp A| = P ∖ |A|), and
% under « Hyp' » (every point is |X| for some X) the Proposition of page 80:
% Σ_M → P(P) is an isomorphism of ordered sets compatible with ∁.
%
% Relation to earlier drafts: GF IV (156-4 batch 4, pp. 65-68) defines
% cosupp, supp, saturated parts and the set Σ of « supports virtuels » in the
% same way, and already presumes an ultra-Stonean Boolean algebra; here the
% Boolean structure is questioned (pp. 63-64) and replaced by realisations
% in P(P). Noted where it bears.
%
% Notation, aligned with 156-6 and 156-4: 𝔉 \mathfrak{F}, 𝓜 \mathcal{M}, ℒ
% \mathcal{L}, 𝔓 \mathfrak{P}; disjointness ∥ \parallel, compatibility
% \between, subdivision \preccurlyeq, incidence \triangleleft; « omb » in
% roman, cosupp and supp as \operatorname; his complement sign ∁ as
% \complement; Σ°_F as \Sigma^{\circ}_{\mathfrak{F}}; the pair of opposed
% arrows between Σ's as \rightleftarrows with the two maps written out.
% Circled labels are given in parentheses with a note.
%
% Hard pages: 63 (NB paragraph dense with interlinear additions), 66
% (second half written between the lines), 70-71 (fast connective prose),
% 73 (NB 2, struck and rewritten), 74 (the example of the string), and the
% oblique left-margin notes on pages 61, 63, 64, 66, 74, 79. Variances kept
% and flagged: « A ∈ F » for A ⊂ F (p. 65); omb_Λ where P is meant (p. 68);
% « A, B » for F, G (p. 68); card X̃ read 0 for a point (p. 70); « X̄ » for
% |X| and « P ∖ A » for P ∖ |A| (p. 77); « = A » for |A| (p. 78); |A| =
% supp A without bars (p. 80).
\documentclass[11pt,a4paper]{article}
\input{../preamble/grothendieck}

\folder{156-7}
\batch{4}
\pages{61}{80}
\dating{1986}
\watermark{Édition de démonstration}
\foldertitle{[Chapitre] VII. Analysis situs (troisième mouture) : notes manuscrites (23-26/06/1986)}

\begin{document}
\page{61}
\note{son numéro « 61 » en haut de la page ; comme dans les chapitres IV et VI, sa pagination, portée en tête de chaque feuillet, coïncide avec celle des archivistes sur tout le lot (pages 61 à 80)}

\subsection{\emph{Notion de support}}
\note{titre souligné, de sa main}

Soit $\mathfrak{F}$ un atelier. Soit $\Sigma_{\mathfrak{F}} \subset \mathfrak{P}(\mathfrak{F})$ l'ens.\ des parties de $\mathfrak{F}$ qui sont \add{(définies par $\leq$)} fermées \struck{à la fois} pour \struck{$\leq$ et} $\ll$, et stables par Sup pour $\leq$ (\uncertain{Sup finis ou infinis}). Soit $\Sigma_{\mathcal{M}}$ l'ens.\ des parties de $\mathcal{M}$ fermées pour $\ll$. Alors on a deux bijections, inverses l'une de l'autre
\[
  (*) \qquad \Sigma_{\mathfrak{F}} \rightleftarrows \Sigma_{\mathcal{M}} , \qquad
  \mathfrak{S} \longmapsto \mathfrak{S} \cap \mathcal{M} , \qquad
  \mathfrak{F}_{S} = \lbrace F \in \mathfrak{F} \mid \widetilde{F} \subset S \rbrace \longleftarrow\!\shortmid\ S
\]
\marginal{\emph{Iso.\ d'ens.\ ordonnés}}
\note{« $(*)$ » : marque lue sans certitude, dans la marge gauche. La page dispose les deux flèches $\rightleftarrows$ entre $\Sigma_{\mathfrak{F}}$ et $\Sigma_{\mathcal{M}}$, avec la première flèche « $\mathfrak{S} \mapsto \mathfrak{S} \cap \mathcal{M}$ » au-dessus (un mot biffé à sa suite) et la seconde au-dessous ; l'indice de $\Sigma_{\mathcal{M}}$ est écrit en surcharge ; « Iso.\ d'ens.\ ordonnés », souligné, est à droite. La lettre rendue $\mathfrak{S}$ est une capitale cursive surchargée}

\uncertain{moyennant} quoi on va les identifier, trouvant un ens.\ $\Sigma$.

Soit \struck{$\mathfrak{F}$} $A \subset \mathfrak{F}$ une partie quelconque. On définit son \emph{cosupport} par
\[
  \operatorname{cosupp} A = \lbrace G \in \mathfrak{F} \mid \forall F \in A,\ F \parallel G \rbrace \subset \mathfrak{F}
\]
\note{le $F$ de « $\forall F \in A$ » est écrit en surcharge}
On voit, grâce à At D, que
\[
  \operatorname{cosupp} A \in \Sigma_{\mathfrak{F}}
\]
\struck{et il est}
On \struck{\ill{}} pose \add{(ou $\overline{A}$)}
\[
  \operatorname{supp} A = \operatorname{cosupp}(\operatorname{cosupp} A) ,
\]
\note{« At D » : lecture douteuse du nom de l'axiome invoqué}
\marginal{Ne pas confondre $\overline{A}$ avec l'adhérence de $A$ pour la topologie de $\mathfrak{F}$ associée soit à $\leq$, soit à $\ll$ !}
\note{note écrite à droite, en regard de « (ou $\overline{A}$) »}

\uncertain{et tout} ceci est formel \uncertain{en termes de la} \uncertain{donnée} d'une relation $\parallel$ dans l'ens.\ abstrait $\mathfrak{F}$ — se \uncertain{généraliserait aussi} pour une relation quelc.\ entre deux ens.)

\marginal{On écrit $F \parallel A$ au lieu de $F \in \operatorname{cosupp} A$. On écrit $A \parallel B$ pour les relations équivalentes :
(i) $A \subset \operatorname{cosupp} B$ ;
(ii) $B \subset \operatorname{cosupp} A$ ;
(iii) $\forall X \in A$, $X \parallel B$ ;
(iii) $\forall Y \in B$, $Y \parallel A$ ;
(iii') $\forall X \in A$, $Y \in B$, $X \parallel Y$ ;
\uncertain{sera} équivalente aussi \uncertain{à} \uncertain{deux autres} \ill{} :
$\operatorname{supp} A \cap \operatorname{supp} B = \emptyset$
(\uncertain{ce} \uncertain{qu'indique} \ill{})}
\note{longue note écrite en oblique dans la marge gauche, sur toute la moitié inférieure de la page. Il numérote deux conditions « (iii) », la première surchargée ; la dernière ligne n'est lue qu'en partie}

\page{62}
\note{son numéro « 62 » en haut de la page}
\struck{\ill{}} $\overline{A}$,

(i) $A \subset B \Longrightarrow \operatorname{cosupp} A \supset \operatorname{cosupp} B$

(ii) $A \subset \overline{A}$, et $A \subset B \Rightarrow \overline{A} \subset \overline{B}$, donc $\bigcup_i \overline{A_i} \subset \overline{\bigcup A_i}$ pour une famille d'ens.\ (mais on n'a pas en général l'égalité)

\struck{(iii)}

(iii) $\overline{\overline{A}} = \overline{A}$

(iv) Conditions équivalentes sur $A \subset \mathfrak{F}$.

\quad a) $\exists\, B$ avec $A = \operatorname{cosupp} B$

\quad b) $A = \overline{A}$

\struck{c)}
On dira alors que $A$ est une partie \emph{\struck{supp}-fermée} de $\mathfrak{F}$, \struck{\ill{}} ou \add{est} un « \struck{\ill{}} $\mathfrak{F}$-\emph{support} ». \struck{\ill{}} On a \uncertain{alors} $A \in \Sigma_{\mathfrak{F}}$.
\note{la parenthèse « (mais on n'a pas en général l'égalité) » est écrite à droite de (ii), sur quatre lignes}

\struck{(v) $\forall A \subset \mathfrak{F}$, $\overline{A}$ est la plus}

Soit $\Sigma^{\circ}_{\mathfrak{F}}$ \struck{$= \ill{}$} la partie de $\Sigma_{\mathfrak{F}}$ formée des \struck{parties fermées de $\mathfrak{F}$. supports de $\mathfrak{F}$.} $\mathfrak{F}$-supports. On a alors
\note{un premier « $\Sigma^{\circ}$ », surchargé, est écrit sous la ligne biffée, avec « supp » au-dessous ; la fin de la ligne suivante est biffée}

(v) Si $A \subset \mathfrak{F}$, $\overline{A}$ est la plus petite \struck{partie} \struck{$\mathfrak{F}$ fermée} $\mathfrak{F}$-support contenant $A$.

(vi) L'application
\[
  A \longmapsto \operatorname{cosupp} A \qquad \mathfrak{P}(\mathfrak{F}) \to \Sigma^{\circ}_{\mathfrak{F}}
\]
\struck{a une application \ill{}} induit une \add{anti}-iso.\ d'ens.\ ordonnés
\[
  \Sigma^{\circ}_{\mathfrak{F}} \longrightarrow \Sigma^{\circ}_{\mathfrak{F}}
\]
qui est \emph{involutif} (exprime le fait que la relation $\parallel$ est symétrique). On pose aussi, si $A$ est un support
\[
  \complement A = \operatorname{cosupp} A ,
\]
\note{« anti » est ajouté au-dessus de la ligne ; devant $\complement A$, un « C » plus appuyé, en marge de la formule}

\page{63}
\note{son numéro « 63 » en haut de la page}
où \uncertain{s'appelle} le \emph{support} \emph{complémentaire} du support donné $A$.

(vii) $\Sigma^{\circ}_{\mathfrak{F}}$ a un plus petit élément (savoir la partie \add{$\emptyset$} \add{réduite à $\emptyset_{\mathfrak{F}}$}$^{*}$ \struck{\ill{}} de $\mathfrak{F}$) et un plus grand élément (savoir $\mathfrak{F}$), \struck{\ill{}} qui sont complémentaires l'un de l'autre.
\note{« $\emptyset$ » est écrit au-dessus de « élément », « réduite à : $\emptyset_{\mathfrak{F}}$ », suivi d'un astérisque, au-dessous de la ligne ; ce qu'il désigne au juste comme plus petit élément reste ainsi indécis sur la page}

\marginal{$^{*}$ Ceci provient du fait que $\emptyset_{\mathfrak{F}}$ est l'unique élément de $\mathfrak{F}$ qui soit $\parallel$ de tout élément de $\mathfrak{F}$ (en \uncertain{nommant} \ill{} $\emptyset_{\mathfrak{F}}$ disjoints d'elles-mêmes…). On le notera, par abus de langage, $\emptyset$ (cf.\ plus haut), \uncertain{ou} $\emptyset_{\Sigma}$}
\note{note écrite en oblique dans la marge gauche, en regard de (vii) et (viii), appelée par l'astérisque}

(viii) \struck{Dans} \add{Plus généralement,} $\Sigma^{\circ}_{\mathfrak{F}}$, il y a des \struck{Sup} \uncertain{Inf} et des \uncertain{Sup} quelconques, \uncertain{échangés} par l'application $\complement$. Le \struck{\ill{}} Inf s'explicite par
\[
  \operatorname{Inf}^{\Sigma^{\circ}_{\mathfrak{F}}}_{i} A_i = \bigcap A_i
\]
et le Sup par
\[
  \operatorname{Sup}^{\Sigma^{\circ}_{\mathfrak{F}}}_{i} A_i = \overline{\textstyle\bigcup A_i} \quad \Longrightarrow \ (= \operatorname{supp}(\textstyle\bigcup A_i))
\]
\note{un signe « $\uparrow \Sigma$ », surchargé, est dans la marge gauche en regard de (ix). La première des deux lectures « Inf » et « Sup » du début de (viii) est écrite sur un mot biffé}

(ix) $\forall A \in \Sigma^{\circ}_{\mathfrak{F}}$, on a
\[
  A \cap \complement A = \emptyset_{\Sigma}
\]
\struck{(En effet \ill{} \ill{} $F$ \ill{})} (Exprime le fait que le seul élément de $\mathfrak{F}$ \struck{\ill{}} tel que $F \parallel F$ est $\emptyset_{\mathfrak{F}}$, lequel est $\parallel$ de $\mathfrak{F}$ tout entier…)

(x) \emph{Cor} $A \subset B \Longrightarrow A \cap \complement B = \emptyset_{\Sigma}$
\note{« (x) » : lecture douteuse du numéro, dans la marge gauche}

\emph{NB} Il n'est pas clair que l'implication $\Rightarrow$ ci-dessus soit \add{(il est \struck{parfois} faux)} \uncertain{claire}. \struck{\ill{}} Il n'est \add{dans $\Sigma^{\circ}_{\mathfrak{F}}$} pas clair que \add{\uncertain{les} Sup} \struck{soit} \add{\uncertain{finis}} \uncertain{soient} \ill{} (contenant que parties de $\mathfrak{F}$), \uncertain{i.e.} dans $\mathfrak{F}$ — la réunion de deux supports, n'est pas \uncertain{a priori} un support : cela exprime le fait \add{\uncertain{intuitif}} qu'une « figure » \add{(continue [\uncertain{dans} \ill{}])} dans la réunion des supports \add{[\ill{} \ill{}]} \ill{} figures n'est pas forcément contenue dans l'un ou l'autre ; elles le \uncertain{seraient} dans le cas des \uncertain{multistrates}…
\note{paragraphe très chargé d'ajouts interlinéaires, dont l'ordre et le rattachement ne sont pas sûrs ; la dernière ligne de la page n'est lue qu'en partie}

\page{64}
\note{son numéro « 64 » en haut de la page}
lieux (quand ce n'est pas dans les ateliers ensemblistes \uncertain{tout au moins}…). Il n'est pas clair non plus, en général, que l'on ait, dans $\Sigma^{\circ}_{\mathfrak{F}}$,
\[
  P \wedge \bigvee_i Q_i = \bigvee_i (P \wedge Q_i)
\]
\note{devant le premier $\bigvee$, un « Sup » est biffé ; devant le second, un signe raturé}
(distributivité de l'intersection des supports, par rapport aux « sommes » même finies…). Pourtant, cela a tendance à être vrai, de même que l'implication \uncertain{réciproque} de $(*)$, dans les \struck{lieux les plus \ill{}} ateliers auxquels j'ai songé jusqu'à présent, cf.\ plus bas.
\note{« $(*)$ » : lecture douteuse du renvoi, qui semble viser l'implication « $\Rightarrow$ » du NB de la page 63}

Toute ceci \uncertain{se} retrouvait dans la \struck{\ill{}} structure $(\mathcal{M}, \leq, \ll, \between)$, \struck{\ill{}} \uncertain{en} \uncertain{codant} \uncertain{dans} en termes de la seule relation $\parallel_{\mathcal{M}}$ \struck{\ill{}} de disjonction dans $\mathcal{M}$. Et l'ens.\ $\Sigma^{\circ}_{\mathcal{M}}$ des « $\mathcal{M}$-supports » \uncertain{est en} bij.\ canonique \ill{} \uncertain{avec} \ill{} \uncertain{canonique}
\[
  \Sigma^{\circ}_{\mathcal{M}} \rightleftarrows \Sigma^{\circ}_{\mathfrak{F}} \qquad \text{iso.\ d'ens.\ ord.}
\]
\note{le signe rendu $\between$ dans le quadruplet est un $\between$ barré, lu sans certitude ; « $\parallel_{\mathcal{M}}$ » a l'indice $\mathcal{M}$ écrit sous le signe}
\uncertain{induite} par les bijections de la p.\ 61) \uncertain{commutant} : l'opération « complémentaire ». Cette fois, le plus petit élément de $\Sigma^{\circ}_{\mathcal{M}}$ est bien la partie \emph{vide} de $\mathcal{M}$, et son plus grand élément est $\mathcal{M}$.

\marginal{Si $A \subset \mathfrak{F}$ est fermé pour $\leq$, \ill{} $\operatorname{cosupp}_{\mathfrak{F}}(A) \cap \mathcal{L} = \operatorname{cosupp}_{\mathcal{L}}(A \cap \mathcal{L})$, en particulier pour $F \in \mathfrak{F}$, $\operatorname{cosupp}_{\mathfrak{F}}(F) \cap \mathcal{L} = \operatorname{cosupp}_{\mathcal{L}}(\ill{}(F))$}
\note{note écrite en oblique dans la marge gauche, au bas de la page ; devant le dernier « $(F)$ », un petit signe surmonté d'un rond n'est pas lu}

\page{65}
\note{son numéro « 65 » en haut de la page}
\struck{On a les \ill{} constructions \uncertain{duales} \uncertain{dans} $\mathcal{L}$ muni de $\parallel$, d'où $\Sigma^{\circ}_{\mathcal{L}}$, ens.\ des parties de $\mathcal{L}$ « fermées ». Considérons les applications canoniques}
\struck{$\Sigma^{\circ}_{\mathcal{M}} \rightleftarrows \overline{\Sigma_{\mathcal{L}}} \overset{\mathrm{déf}}{=} \mathfrak{P}(\mathcal{L})$, $A \mapsto A \cap \mathcal{L}$, $A_S = \lbrace X \in \mathcal{M} \mid \mathrm{omb}(X) \subset S \rbrace \longleftarrow\!\shortmid\ S$}
\note{ces lignes, jusqu'à la formule, sont barrées de trois longs traits obliques ; le « $\overline{\phantom{x}}$ » sur $\Sigma_{\mathcal{L}}$ est un trait appuyé qu'on lit sans certitude comme une barre}

Cette relation entre supports dans $\mathfrak{F}$, et supports dans $\mathcal{M}$, \uncertain{se} \uncertain{généralisait} en remplaçant $\mathcal{M}$ par une sous-ensemble $\Lambda \subset \mathfrak{F}$, soumis à la seule condition que voici :

\emph{Disj}$_{\Lambda}$ : Soit $F \in \mathfrak{F}$, $\mathrm{omb}_{\Lambda}(F) = \lbrace X \in \Lambda \mid X \ll F \rbrace$. Alors pour tout $G \in \mathfrak{F}$, on a
\[
  F \parallel G \iff \forall X \in \mathrm{omb}_{\Lambda}(F),\ X \parallel G
\]
\note{l'énoncé de Disj$_{\Lambda}$ est marqué d'un trait vertical dans la marge gauche ; « Disj » est souligné}

Posant, pour $A \subset \mathfrak{F}$
\[
  \mathrm{omb}_{\Lambda}(A) \overset{\mathrm{déf}}{=} \bigcup_{F \in A} \mathrm{omb}_{\Lambda}(F) \subset \Lambda ,
\]
on voit donc que pour toute partie $A \in \mathfrak{F}$, on a
\[
  \operatorname{cosupp}_{\mathfrak{F}} A = \operatorname{cosupp}_{\mathfrak{F}} \mathrm{omb}_{\Lambda} A .
\]
\note{« $A \in \mathfrak{F}$ » est ainsi sur la page, pour $A \subset \mathfrak{F}$}
De plus, si $A$ est fermé dans \struck{$\mathfrak{F}$} $\mathfrak{F}$ pour $\ll$, alors on voit aussitôt que
\[
  \mathrm{omb}_{\Lambda}(A) = \Lambda \cap A \qquad (A \text{ fermé dans } \mathfrak{F}, \ll)
\]
(en général on a seulement $\supset$). Donc on a \uncertain{cependant}
\[
  \operatorname{cosupp}_{\mathfrak{F}}(A) = \operatorname{cosupp}_{\mathfrak{F}}(A \cap \Lambda) \qquad \text{si } A \text{ fermé dans } \mathfrak{F} \text{ pour } \ll
\]
Ceci implique, en reprenant les cosupports
\struck{$A$}
\[
  \operatorname{supp}_{\mathfrak{F}}(A) = \operatorname{supp}_{\mathfrak{F}}(A \cap \Lambda) \qquad \text{si } A \text{ fermé dans } \mathfrak{F} \text{ pour } \ll
\]
et \uncertain{sans hyp.\ sur} $A$
\[
  \operatorname{supp}_{\mathfrak{F}}(A) = \operatorname{supp}_{\mathfrak{F}}(\mathrm{omb}_{\Lambda}(A))
\]
\note{dans la formule $\operatorname{supp}_{\mathfrak{F}}(A) = \ldots$, le premier « $\operatorname{supp}$ » est écrit en surcharge sur un « $A$ » raturé, et le second sur une lettre surchargée ; « $A \cap \Lambda$ » est placé en exposant au-dessus de « $\operatorname{cosupp}_{\mathfrak{F}}$ » dans les deux formules}

\page{66}
\note{son numéro « 66 » en haut de la page}
On trouve ainsi une correspondance biunivoque entre les ens.\ de $\mathfrak{F}$-supports et de $\Lambda$-supports
\[
  \Sigma^{\circ}_{\mathfrak{F}} \rightleftarrows \Sigma^{\circ}_{\Lambda} , \qquad
  A \longmapsto A \cap \Lambda = \mathrm{omb}_{\Lambda}(A) , \qquad
  \operatorname{supp}_{\mathfrak{F}} S = \lbrace F \in \mathfrak{F} \mid \mathrm{omb}_{\Lambda}(F) \subset S \rbrace \longleftarrow\!\shortmid\ S
\]
\note{la page dispose les deux flèches entre $\Sigma^{\circ}_{\mathfrak{F}}$ et $\Sigma^{\circ}_{\Lambda}$, puis au-dessous les deux applications ; « $\operatorname{supp}_{\mathfrak{F}} S$ » est relié par un signe $=$ vertical à l'ensemble $\lbrace F \in \mathfrak{F} \mid \ldots \rbrace$ écrit sous lui}
respectant les relations d'ordre (inclusion) et la « \uncertain{passage} au complémentaire », que je vais noter $\complement_{\mathfrak{F}}$, $\complement_{\Lambda}$.
\note{« passage » : le mot est écrit « p.r.g. » ou approchant ; on lit la même abréviation à la page 68}

On a intérêt, pour \add{« comprendre » les supports}, de prendre $\Lambda$ aussi petit que possible, \uncertain{mais cela demande} \ill{}. Le choix $\mathcal{L}$ s'impose, mais cela demande que l'axiome At L 2 soit satisfait — ce qui n'est pas \uncertain{toujours} \uncertain{cependant}. Ainsi, on trouve (\uncertain{dans ce cas}) \uncertain{par} \uncertain{abus}, \uncertain{en} \ill{} \uncertain{laissant} tomber l'exposant \ill{} \ill{} \ill{} \ill{} $\Sigma$…)
\[
  \Sigma_{\mathfrak{F}} \simeq \Sigma_{\mathcal{M}} \simeq \Sigma_{\mathcal{L}}
\]
\ill{} \ill{} sous-entendu : chaque fois : la notation d'\uncertain{iso.} \uncertain{ensemblistes}, \uncertain{des} supports, \uncertain{définie} \uncertain{ensemblistement}, en \uncertain{termes} \uncertain{de} $\parallel_{\mathfrak{F}}$ \uncertain{ou} $\parallel_{\mathcal{M}}$ \uncertain{ou} $\parallel_{\mathcal{L}}$).
\note{« At L 2 » est un nouveau nom d'axiome, pour les lieux ; il revient aux pages 68 et 73. Le passage qui suit est écrit vite, entre les lignes, et n'est lu qu'en partie}

\marginal{\emph{NB} Les \ill{} \ill{} \ill{} \ill{} de $\mathfrak{F}$ \ill{} \ill{} \ill{} At L 2, il \ill{} \ill{} \ill{} $\mathcal{M}$ \ill{} \ill{} \ill{} \ill{} de $\mathcal{L}$, \ill{} \ill{} \ill{} de la relation $\parallel_{\mathcal{L}}$}
\note{note écrite en oblique dans la marge gauche, au bas de la page, lue par fragments}

Un cas particulièrement agréable est celui où $\parallel_{\mathcal{L}}$ est la relation tautologique : $x \neq y$ (p.ex.\ cas d'un atelier quasi-ensembliste)
\note{« agréable » est souligné}

\page{67}
\note{son numéro « 67 » en haut de la page}
on trouve
\[
  \Sigma_{\mathcal{L}} = \mathfrak{P}(\mathcal{L}) ,
\]
l'opération $\complement_{\mathcal{L}}$ étant simplement le passage au complémentaire. Dans ce cas, le support d'une partie $A$ de $\mathfrak{F}$ (ou de $\mathcal{M}$, \uncertain{quand on y passe en strates}), identifié à une partie de $\mathcal{L}$, n'est autre que $\mathrm{omb}_{\mathcal{L}}(A) = \mathrm{omb}(A)$. Donc dans ce cas, la définition devient simplement : posons
\[
  \operatorname{supp} A = \mathrm{omb}(A) \overset{\mathrm{déf}}{=} \bigcup_{F \in A} \mathrm{omb}(F)
\]
\[
  \operatorname{cosupp} A = \complement_{\mathcal{L}} \operatorname{supp} A \qquad \text{(complémentaire ordinaire)}
\]
Ceci traduit en fait que la relation cruciale de disjonction, devient simplement
\[
  F \parallel G \iff (\mathrm{omb}\, F \cap \mathrm{omb}\, G = \emptyset)
\]
relation de disjonction ensembliste sur les ombres (relation \uncertain{d'objets} supports pour la \uncertain{omb} \uncertain{ensembliste}). Dans ce cas, les perplexités p.\ 63, 64 se résolvent, on a
\[
  \left\lbrace
  \begin{array}{l}
    A \subset B \iff A \cap \complement(B) = \emptyset_{\Sigma} \\
    A \cap \bigvee_i B_i = \bigvee_i A \cap B_i \quad \text{dans } \Sigma \ —
  \end{array}
  \right.
\]
\note{une flèche relie le $\emptyset$ de la formule de $F \parallel G$ au mot « ensembliste » de la ligne suivante, dont la fin n'est lue qu'en partie ; sous $\complement(B)$, une petite flèche montante et le mot « compl. » entre guillemets}

\page{68}
\note{son numéro « 68 » en haut de la page. Un long trait oblique traverse toute la page, de haut en bas ; on transcrit le texte, qui se lit dessous}
Ceci peut se généraliser ainsi. Soit
\[
  P \subset \mathcal{L}
\]
\marginal{$P$ comme « points »}
une partie de $\mathcal{L}$ (ou simplement de $\mathfrak{F}$ — mais pratiquement \uncertain{on sera} \uncertain{toujours} $\subset \mathcal{L}$ je présume) telle que

a) si $x, y \in P$, $x \neq y \Rightarrow x \parallel y$ (donc sur $P$, $x \parallel y$ équivaut à $x \neq y$)

b) si $F, G \in \mathfrak{F}$, (\uncertain{resp.} $\mathcal{M}$, \ill{} \ill{} \ill{} $\widetilde{\mathcal{M}}$) \struck{alors et si $\forall x \in \mathrm{omb}_{\Lambda}(F)$, $y \in \mathrm{omb}_{\Lambda}(G)$, \ill{}} \struck{\ill{}} on a $F \parallel G$ si (et ss.) $\nexists\, x \in P$, \struck{\ill{}} avec $\operatorname{supp} x \subset \operatorname{supp} A \cap \operatorname{supp} B$ \struck{\ill{} $x \parallel y$, \ill{}} (\struck{C'est la condition \ill{} At $\Lambda$ 2 dans le cas $\Lambda = P$} $=$ condition qui \uncertain{remplace} At L 2.)
\note{b) est surchargé de ratures : trois lignes et demie sont biffées d'un trait horizontal et de traits obliques, et l'on ne donne que ce qui se lit ; « (et ss.) » est écrit au-dessus d'un mot raturé. Les lettres $A$, $B$ de la dernière formule sont sur la page, là où l'on attend $F$, $G$}

\marginal{Un peu bref !}
\note{écrit en oblique dans la marge gauche, avec une flèche descendant vers « Alors »}

Alors l'ens.\ ordonné des supports s'identifie à
\[
  \Sigma_{P} \simeq \mathfrak{P}(P) , \qquad \text{donc} \quad A \longmapsto \lbrace x \in P \mid \operatorname{supp} x \subset A \rbrace
\]
\[
  \Sigma_{\mathfrak{F}\ (\text{ou } \mathcal{M},\ \text{ou } \mathcal{L})} \simeq \mathfrak{P}(P)
\]
\note{dans la première formule, la page écrit « $\Sigma_{P} \simeq \mathfrak{P}(P)$, donc », puis au-dessous « $A \mapsto \lbrace x \in P \mid \operatorname{supp} x \subset A \rbrace \simeq \mathfrak{P}(P)$ » ; l'indice $\mathfrak{F}$ de la seconde est surchargé}
iso d'ens.\ ordonnés, compatible à l'application $\complement$ (« complémentaire »). Donc dans ce cas aussi, les perplexités se dissipent. Dans ce cas, il y a lieu de définir le support comme
\[
  \operatorname{supp} A = \mathrm{omb}_{\Lambda} A = \bigcup_{F \in A} \mathrm{omb}_{\Lambda}(F) .
\]
Notons que l'isom.\ d'ens.\ ordonnés avec application « passage au complémentaire »
\[
  \Sigma \simeq \mathfrak{P}(P)
\]
\note{dans les deux dernières formules l'indice de $\mathrm{omb}$ est un $\Lambda$ ; on attend $P$}

\page{69}
\note{son numéro « 69 » en haut de la page}
\subsection{25 juin}
\note{date de sa main, « (25 juin) », entre parenthèses, dans l'angle supérieur gauche, avec un trait courbe qui la rattache au début du texte}

En dehors du cas quasi-ensembliste, je suspecte que la notion de support, telle qu'elle est présentée ici, est caduque — qu'il ne faudrait parler de support d'une partie $A$ de $\mathcal{M}$, disons, que quand les $X \in A$ \struck{étant} sont mutuellement compatibles, ou \ill{} \uncertain{unions}, si deux à deux ils admettent des \ill{} mutuellement compatibles. Voici \uncertain{une situation} \ill{} un exemple typique.

\emph{Exemple} \quad Strates associées à un \struck{segment} \add{ens.} ordonné $I$ (p.ex.\ un « \emph{segment} ordonné »).
\[
  \mathcal{M} = \mathcal{M}_0 \amalg \mathcal{M}_1
\]
$\mathcal{M}_0$ (0-strates) $= \struck{\ill{}}\ (= \mathrm{Drap}_0(I))$,

$\mathcal{M}_1$ (1-strates) $=$ ens.\ des sous-segments de $I$, i.e.\ parties \add{$X$} de la forme $I_{a \leq \cdot \leq b}$, où $a < b$. (NB $a$, $b$ sont uniquement déterminés : origine et extrémité de $X$. On pose $\partial X = \lbrace a, b \rbrace$.
\note{« Drap » : lecture douteuse. Une première rédaction des lignes de $\mathcal{M}_1$, encadrée puis biffée, se lit en partie : « $\mathcal{M}_1$ (1-strates) $= \mathrm{Drap}_1(I) \simeq \lbrace (a,b) \in I \times I \mid a < b \rbrace$ », « ens.\ des parties \ldots{} de $I$ de la forme $X_{ab} = \lbrace x \mid a \leq x \leq b \rbrace$ où $a < b$. Alors $a$, $b$ uniquement déterminés, on écrira $\partial X = \lbrace a, b \rbrace \subset \mathcal{M}_0$ », et dans le cadre : « $X \leq Y$ ssi $X \subset Y$ dans \ldots{} $X < Y$ ssi $X \in \mathcal{M}_1$, $Y \in \mathcal{M}_1$, $X \in \partial Y$ ; $X \ll Y$ ssi \ldots{} »}
\[
  X \leq Y \iff \left\lbrace
  \begin{array}{l}
    X = Y \ \text{ou} \\
    X \in \mathcal{M}_0,\ Y \in \mathcal{M}_1,\ X \in \partial Y
  \end{array}
  \right.
\]
\marginal{\struck{Cela \ill{} \ill{} \ill{} $X \mathrel{\triangleleft} Y$} \ill{} $X \mathrel{\triangleleft} Y$}
\note{note à droite de la formule, où « $X \mathrel{\triangleleft} Y$ » est écrit deux fois, la première dans une ligne biffée ; dans « $X \in \partial Y$ », le $Y$ est écrit en surcharge}
\[
  X \ll Y \iff X \subset Y
\]
\[
  X \between Y \iff \text{on est dans un des trois cas mutuellement exclusifs}
\]
\note{« mutuellement exclusifs » est lu dans la parenthèse « ($X \overset{\circ}{\ll} Y$ signifie \ill{} \ldots{} mutuellement exclusifs) » écrite au-dessus, en partie biffée ; le « $\overset{\circ}{\ll}$ » y est lu sans certitude}

(a) $X = Y$ \quad (b) $X \in \partial Y$ ou $Y \in \partial X$ (cas où $Y$ majore $X$ et dans $\mathcal{M}_1$). \quad (c) $\partial X \cup \partial Y$ tot.\ ordonné, et $X \cap Y$ est vide, ou réduit à un point, \uncertain{dire} \struck{\ill{}} ! \uncertain{un point}, \ill{} \ill{} \uncertain{écrit dans un des cas} : ($\alpha$) $\exists\, a < b < c < d$ avec $X = [a,b]$, $Y = [c,d]$, ($\beta$) $\exists\, a < b < c$ avec $X = [a,b]$, $Y = [b,c]$, ou l'inverse.
\note{« (a) », « (b) », « (c) » sont cerclés ; dans (c), plusieurs mots à la suite de « tot.\ ordonné » sont raturés à droite. À gauche, trois petits croquis : pour b), un segment $X$ et un point $Y$ à son extrémité ; pour c), deux segments $X$, $Y$ séparés, « ou » deux segments $X$, $Y$ bout à bout}

\emph{NB}. Tous les axiomes sont satisfaits, à l'exception de At L 1 (\uncertain{admissibilité}), satisfait ssi les intervalles $]a,b[$ sont $\neq \emptyset$ ($a < b$).

\page{70}
\note{son numéro « 70 » en haut de la page}
\emph{NB} Les strates sont combinatoirement finies — celles de $\mathcal{L} = \mathcal{M}_0$ sont de dim $0$, celles de $\mathcal{M}_1$ de dim $1$. Si $X \in \mathcal{M}_0$, \struck{$X$} $\operatorname{card} \widetilde{X} = \uncertain{0}$, si $X \in \mathcal{M}_1$, $\operatorname{card} \widetilde{X} = 3$. On prend pour $\mathfrak{F}$ les figures \emph{finies}, i.e.\ les parties finies $F$ de $\mathcal{M}$, \struck{telles que} fermées pour $\leq$ (i.e.\ contenant, avec $X \in \mathcal{M}_1$, ses deux extrémités) et telles que $X, Y \in F \Rightarrow X \between Y$.
\note{« finies » est souligné deux fois. Le chiffre lu $0$ pour $\operatorname{card} \widetilde{X}$, $X \in \mathcal{M}_0$, est petit et peu formé ; avec $\widetilde{X} = \lbrace Y \mid Y \leq X \rbrace$ on attendrait $1$, comme le $3$ du cas $X \in \mathcal{M}_1$ compte $X$ et ses deux extrémités}

\struck{Comme \ill{} dans GF III j'ai appelé « parties}

\drawing{sous cette ligne biffée, un croquis : une droite portant, de gauche à droite, un segment subdivisé en deux, un point isolé, un segment, deux points, un segment plus long subdivisé, un segment finement subdivisé, un point — une figure finie de l'exemple}

Dans GF III, j'ai pris la peine essentiellement de définir le support d'un tel objet comme partie de $I$ (appelée \struck{\ill{}} « modérée fermée » \uncertain{dont} j'aurais dû dire « modérée » tout court, et appeler \emph{constructibles} les parties que j'appelle « modérées »). Cela donne lieu à un formalisme raisonnable tant \uncertain{que} la \uncertain{finitude} \uncertain{d'un} ens.\ de figures \uncertain{sujet à} la propriété spécifiée au début. \uncertain{Dans} \uncertain{le} \uncertain{cas} \uncertain{d'un} \ill{}, c'est OK si $I$ est un ensemble \ill{}, \uncertain{c'est} \ill{} (loc.\ cit.\ p.\ 8) et \uncertain{aucun} \emph{tronçon} ordonné. Toujours est-il que l'on trouve une notion de parties \emph{constructibles} de $I$, de complémentaire d'une telle partie, \struck{et par} une notion de compatibilité de parties constructibles, et une bonne formalisme de $\cup$, $\cap$, $\complement$ pour les parties constructibles \add{$X_\alpha$} deux à deux \struck{disjointes} compatibles (i.e.\ telles que la réunion des $\partial X_\alpha$ soit tot.\ ordonnée).
\note{« GF III » est le chapitre III, dossier 156-3 ; « loc.\ cit.\ p.\ 8 » y renvoie. La phrase du milieu, écrite vite, n'est lue qu'en partie}

\page{71}
\note{son numéro « 71 » en haut de la page}
Je pense qu'on \uncertain{aura} un tel formalisme sans hyp.\ sur $I$ (genre « \uncertain{tronçons} ordonnés ») à condition de tenir compte de la possibilité qu'il y ait des 1-strates
\[
  X = [a,b] \qquad (a < b)
\]
réduites à $\lbrace a, b \rbrace$ (quand $]a,b[ = \emptyset$). Dans ce cas, il serait déraisonnable de définir $\operatorname{supp} X = [a,b]$ ou, donc $= \struck{\ill{}}\ \operatorname{supp} \partial X$. Mais on doit pouvoir définir une notion de supports, \uncertain{associée} à une partie tot.\ ordonnée finie quelconque $\Delta$ de $I$ (à savoir, $\Delta = \bigcup_\alpha \partial X_\alpha$)
\[
  a_1 < a_2 < \cdots < a_n
\]
comme étant \struck{\ill{}} les éléments de l'algèbre de Boole sur l'ens.\ des « points »
\[
  \lbrace a_1 \rbrace ,\ [a_1, a_2] ,\ \lbrace a_2 \rbrace ,\ [a_2, a_3] ,\ \ldots ,\ \lbrace a_n \rbrace
\]
\struck{Pour} Toute \struck{famille} $\mathfrak{X} \subset \mathcal{M}$, $\mathfrak{X} = \lbrace X_\alpha \rbrace$, \emph{subordonnée} à $\Delta$ i.e.\ telle que $\partial X_\alpha \subset \Delta$ $\forall \alpha$, on définit $\operatorname{supp} \mathfrak{X}$ comme partie de cet ens.\ (partie \uncertain{ens.\ ordonné} fermée, en fait, \ill{} \ill{} \ill{} \ill{} l'inclusion \ill{} \ill{} de $I$), et on définit de \uncertain{même} \ill{} \uncertain{supports}, d'une \uncertain{définition} \uncertain{virtuelle} $X \smallsetminus Y$ pour $Y \ll X$, d'une \uncertain{définition} $F \smallsetminus G$ pour $F \leq G$, $F$, $G$ figures dans la \struck{\ill{}} partie $\leq$-fermée de $\mathfrak{F}$ engendrée par $\mathfrak{X}$.
\note{« subordonnée » est souligné. La lettre rendue $\mathfrak{X}$ est une capitale cursive, la même que pour $\mathfrak{X} = \lbrace X_\alpha \rbrace$ ; « $\Delta$ » est tantôt un triangle plein, tantôt un triangle ouvert. La parenthèse sur la partie fermée, écrite serré avec des ajouts, n'est lue qu'en partie}

Mais \uncertain{ceci} signifie \uncertain{justement} que pour \struck{\ill{}} les figures dans un atelier, i.e.\ \quad \ldots )
\note{la phrase s'achève sur un croquis : un segment subdivisé, dont les points marqués sont étiquetés $a_1$, $a_2$, \ldots, $a_n$, suivi d'une parenthèse fermante}

\page{72}
\note{son numéro « 72 » en haut de la page}
on a une notion de « supports \uncertain{modérés} » \struck{\ill{}} \uncertain{associés} à cette figure, l'ens.\ \add{\uncertain{ordonné}} de ses supports s'identifiant à $\mathfrak{P}(\widetilde{F})$. Et si $F' \preccurlyeq F$ est une subdivision, ou \uncertain{plus} \uncertain{finement} $F' \ll F$ un raffinement, on \uncertain{aura} \uncertain{correspondance} de l'un à l'autre, via
\[
  \varphi : \widetilde{F'} \longrightarrow \widetilde{F}
\]
\uncertain{induisant}
\[
  \varphi^{*} : \mathfrak{P}(\widetilde{F}) \longrightarrow \mathfrak{P}(\widetilde{F'}) .
\]
\note{au-dessus de « supports », un mot interlinéaire est lu « ordo… » sans certitude ; le signe de « $F' \preccurlyeq F$ » est son $\leq$ à trait inférieur courbé}
Je ne \uncertain{soupçonnerais} pas, dans le cas d'un atelier qui n'est pas quasi-ensembliste, qu'on puisse interpréter \struck{\ill{}} de façon canonique les $\mathfrak{P}(\widetilde{F})$, pour des figures différentes (et qui ne seraient ni compatibles, ni reliées par $F' \ll F$, ni \add{pré}compatibles dans le sens qu'elles deviennent compatibles après subd.\ de l'une et de l'autre…), comme des parties d'un même
\[
  \mathfrak{P}(P) \quad \text{disons, où } P \text{ est}
\]
un ensemble de « points » associé d'une façon ou d'une autre au \struck{\ill{}} strates $\mathcal{M}$.

On va définir cependant la notion de « réalisation par supports » d'un atelier, dans un ensemble $P$ de « points ».

\page{73}
\note{son numéro « 73 » en haut de la page}
C'est la donnée d'une application
\[
  \mathcal{M} \longrightarrow \mathfrak{P}(P) \qquad X \longmapsto |X|
\]
satisfaisant les conditions

(i)$_{\mathcal{M}}$ \quad $X \ll Y \Longrightarrow |X| \subset |Y|$ \quad \add{application \emph{croissante} (pour $\ll$, $\subset$)}

(ii)$_{\mathcal{M}}$ \quad $X \between Y \Longrightarrow |X| \cap |Y| = \bigcup_{Z \in X \cap Y} |Z|$

(iii)$_{\mathcal{M}}$ \quad $X \parallel Y \iff |X| \cap |Y| = \emptyset$ \quad (NB $\Rightarrow$ \uncertain{par} (ii))

(iv)$_{\mathcal{M}}$ \quad \struck{\ill{}} $X \preccurlyeq Y \Longrightarrow |X| = |Y|$

(v)$_{\mathcal{M}}$ \quad $|X| = \bigcup_{Z \in \mathrm{omb}(X)} |Z|$

(vi)$_{\mathcal{M}}$ \quad (b) si $x, y \in P$, $\exists\, X \in \mathcal{M}$ avec $x \in |X|$, $y \notin |X|$ ou l'inverse \quad (a) $\bigcup_{X \in \mathcal{M}} |X| = P$, et
\note{le « (iii)$_{\mathcal{M}}$ » est entouré d'un cercle, et une flèche y mène depuis la marge ; le numéro de (vi) est écrit en surcharge sur un autre ; « (a) » et « (b) » sont cerclés, (a) est entouré d'un trait qui le renvoie par une flèche avant (b) — il faut sans doute lire (a) avant (b). Dans (ii), l'indice de la réunion, écrit sous la formule, est lu « $Z \in X \cap Y$ » sans certitude ; dans (i), l'ajout « $|Z|$ » au-dessus de la ligne n'est pas rattaché}
\marginal{le fond de l'axiome est dans l'implication $\Leftarrow$}
\note{écrit dans la marge gauche, avec une flèche vers (iii)$_{\mathcal{M}}$}

\emph{NB} 1) \struck{\ill{}} (L'axiome (vi) est dans la nature d'une normalisation ; il suffit de remplacer $P$ par son quotient par la relation d'équivalence
\[
  x \sim y \iff \mathcal{M}_x = \mathcal{M}_y \qquad (\text{où } \mathcal{M}_x = \lbrace X \in \mathcal{M} \mid x \in |X| \rbrace)
\]

\emph{NB} 2) La condition v implique qu'il suffit de connaître $|Z|$ pour $Z \in \mathcal{L}$. \struck{Mais cela} \struck{\ill{}} \uncertain{On} \emph{définit} \uncertain{alors} $|X|$ pour $X$ général. Mais la validité des conditions (i) : (iv) \uncertain{dépend} des conditions \uncertain{sur} les $|Z|$. La condition (iii) p.ex.\ équivaut \uncertain{à} \add{(ou (iii'))} \uncertain{pour} $x, y \in \mathcal{L}$, $+$ la condition \struck{\ill{}} \add{(ii)}, $\mathcal{L}$ $=$ condition \uncertain{(vi')} \struck{\ill{} \ill{}} At L 3 et (iii') At L 2. \struck{La condition (iii') : \ill{} $X_i$} La condition (i) tautologique (\uncertain{variante} $\Rightarrow$). La condition (iv) \uncertain{résulte} des \uncertain{autres} \add{(moyennant At L 2, L 3)} \uncertain{via} \ill{} \ill{} de \uncertain{supposer} \uncertain{ces} conditions. Donc, la variante
\[
  \mathcal{L} \longrightarrow \mathfrak{P}(P) , \qquad x \longmapsto |x| ,
\]
\note{le second NB, écrit vite et chargé de ratures et d'ajouts, n'est lu qu'en partie ; « At L 3 » y apparaît pour la première fois dans ce lot. La phrase se poursuit page 74}

\page{74}
\note{son numéro « 74 » en haut de la page}
conditions

(i)$_{\mathcal{L}}$ \quad Soient $x, y \in \mathcal{L}$, alors \add{$x \neq y$} $x \parallel y \iff |x| \cap |y| = \emptyset$

(ii)$_{\mathcal{L}}$ \quad \struck{\ill{}} $\forall\, \xi, \eta \in P$, \add{$\xi \neq \eta$}, $\exists\, x \in \mathcal{L}$, avec $\xi \in |x|$, $\eta \notin |x|$ \uncertain{ou l'inverse}
\note{les deux conditions sont marquées d'un trait vertical dans la marge ; dans (i)$_{\mathcal{L}}$, un signe surchargé précède « $x \parallel y$ » ; « ou l'inverse », écrit sous la fin de (ii) et cerclé d'un trait, est suivi d'un mot biffé}

\emph{Exemple} \quad $P$ est une ficelle topologique, et $\mathcal{M}$ les strates des \add{(lieux et des)} segments (ou bouts-ficelles), la fonction support $|x|$ \struck{\ill{}} du \uncertain{dit}, \uncertain{la} \uncertain{fonction} \uncertain{pour} les lieux étant la \uncertain{fonction} \uncertain{évidente}. La condition (ii) n'est pas \uncertain{satisfaite} (\ill{} \ill{} \ill{}, si les bouts de \ill{} \ill{} \ill{} : des points de la ficelle ne \uncertain{sont} \ill{} \ill{} \ill{}, ou \ill{} \ill{} \ill{} \ill{} \ill{}…) \uncertain{mais} (vi) \uncertain{sont} \uncertain{satisfaites}…)
\note{« Exemple » est souligné. La fin de l'exemple, écrite vite, n'est lue qu'en partie}

\marginal{\emph{NB} \ill{} \ill{} pour la \uncertain{condition} (ii), \uncertain{que} les « \ill{} » \ill{} \ill{} \uncertain{ficelles} \uncertain{communes}}
\note{note écrite en oblique dans la marge gauche, en regard de l'exemple, lue par fragments}

\emph{Dans} tous les cas \uncertain{intéressants} de strates \uncertain{connus}, \uncertain{sauf} \uncertain{peut-être} \uncertain{celui} des strates associé à un ens.\ ordonné (\uncertain{traversé} \ill{} \ill{}) donné \uncertain{ad hoc}, on \uncertain{trouve} une notion naturelle de supports $|\ |$.

La question qui se pose donc aussitôt : Une telle donnée détermine-t-elle $P$ à isomorphisme unique près — p.ex.\ par l'ex\struck{istence} d'un iso com.\ \add{d'ens.\ ordonnés avec} \struck{application} \add{opération} $\complement$
\[
  \Sigma_{\mathcal{M}} \simeq \mathfrak{P}(P) \ ??
\]
\note{sous $\Sigma_{\mathcal{M}}$, « ens.\ des supports “abstraits” de $\mathcal{M}$ »}

\page{75}
\note{son numéro « 75 » en haut de la page}
Pour toute partie $A$ de $\mathcal{M}$, posons
\[
  |A| = \bigcup_{X \in A} |X|
\]
\marginal{\emph{NB} $|\emptyset| = \emptyset$, $\ldots$ \uncertain{vide} \uncertain{ainsi}\ldots{} \uncertain{si} $|\mathcal{M}| = P$}
\note{note écrite en oblique dans la marge gauche, en regard de la formule ; le dernier $|\mathcal{M}|$ est surchargé}

Ceci posé, je \add{voudrais} \struck{dis} que pour deux parties $A$, $B$ de $\mathcal{M}$, on ait
\[
  (*) \qquad \operatorname{supp} A \subset \operatorname{supp} B \overset{?}{\iff} |A| \subset |B|
\]
(ou encore $A \subset \operatorname{supp} B$). En effet, \struck{on} la première \struck{\ill{}} \add{\uncertain{inclusion}} équivaut à
\[
  (\alpha) \qquad \operatorname{cosupp} B \subset \operatorname{cosupp} A ,
\]
\note{le point d'interrogation est au-dessus du $\iff$ ; « (ou encore $A \subset \operatorname{supp} B$) » est écrit sous le membre de gauche ; « ($\alpha$) » est surchargé}
or \struck{on a} on a, pour $X \in \mathcal{M}$
\[
  \begin{array}{l}
  X \in \operatorname{cosupp} A \overset{\mathrm{déf}}{\iff} X \parallel A \overset{\mathrm{déf}}{\iff} \forall Y \in A,\ X \parallel Y \\
  \qquad \overset{\mathrm{iii}_{\mathcal{M}}}{\iff} \forall Y \in A,\ |X| \cap |Y| = \emptyset \iff |X| \cap |A| = \emptyset
  \end{array}
\]
Donc l'inclusion $\operatorname{supp} A \subset \operatorname{supp} B$, ou encore $(\alpha)$, équivaut à
\[
  (\cdot) \qquad \forall X \in \mathcal{M}, \quad |X| \cap |B| = \emptyset \Longrightarrow |X| \cap |A| = \emptyset .
\]
et cette \add{propriété} \struck{\ill{}} sur $(A, B)$ est évidemment impliquée par $|A| \subset |B|$. Donc dans (*) on a $\Leftarrow$. Pour voir $\Rightarrow$, il revient au même de montrer que \uncertain{dès} $|A| \not\subset |B|$ implique l'existence de $X \in \mathcal{M}$, avec $|X| \cap |B| = \emptyset$ mais $|X| \cap |A| \neq \emptyset$. Ce n'est pas clair en général, \uncertain{et} je \uncertain{donnerai} plus \uncertain{loin} \uncertain{un} contre-exemple \uncertain{typique}. Mais supposons que $B$ soit un support, i.e.\ \struck{$B$}
\[
  B = \operatorname{supp} B ,
\]
\note{la marque de la condition « $(\cdot)$ » est lue sans certitude}

\page{76}
\note{son numéro « 76 » en haut de la page}
alors le 1er membre de (*) \struck{\ill{}} équivaut à
\[
  A \subset B
\]
qui implique bien $|A| \subset |B|$ \uncertain{tautologiquement}. Ainsi

\emph{Lemme} ! Sous la \add{seule} condition (iii$_{\mathcal{M}}$), si $A$ et $B$ sont deux parties de $\mathcal{M}$ on a
\[
  \operatorname{supp} A \subset \operatorname{supp} B \Longleftarrow |A| \subset |B|
\]
\[
  (\Leftrightarrow A \subset \operatorname{supp} B)
\]
et l'implication inverse est vraie si $B$ est un support (i.e.\ $B = \operatorname{supp} B$).
\note{« Lemme ! » est souligné. Sur la page, « $(A \subset \operatorname{supp} B)$ » est écrit sous le membre de gauche, relié à lui par un signe $=$ vertical ; la parenthèse est rendue ici par $\Leftrightarrow$}

\emph{Corollaire} \quad L'application
\[
  \Sigma_{\mathcal{M}} \longrightarrow \mathfrak{P}(P) \qquad A \longmapsto |A|
\]
est injective, et la relation d'inclusion dans $\Sigma_{\mathcal{M}}$ est induite par la relation d'inclusion dans $\mathfrak{P}(P)$.
\note{« Corollaire » est souligné ; son énoncé est marqué d'un trait vertical dans la marge}

Qu'en est-il \struck{de} de l'intersection $\bigwedge_i$ \add{$= \bigcap_i$} et du \struck{la « réunion »} sup $\bigvee_i$ des supports ? On a
\[
  \left| \bigwedge_i A_i \right| \subset \bigcap_i |A_i| \qquad \text{tautologiques}
\]
\[
  \left| \bigvee_i A_i \right| \supset \bigcup_i |A_i| \qquad \text{id}
\]
mais je doute qu'on ait tjrs égalité. D'autre part, je \uncertain{dis} que
\[
  |\complement_{\mathcal{M}}(A)| \subset P \smallsetminus |A|
\]
\note{dans la dernière formule, le $P$ est écrit en surcharge sur une autre lettre ; la page s'arrête sur cette formule}

\page{77}
\note{son numéro « 77 » en haut de la page}
En effet,
\[
  \complement_{\mathcal{M}}(A) = \lbrace X \in \mathcal{M} \mid X \parallel A \rbrace = \lbrace X \in \mathcal{M} \mid |X| \cap |A| = \emptyset \rbrace
\]
\note{sous le dernier membre : « i.e.\ $|X| \subset P \smallsetminus A$ »}
donc
\[
  |\complement_{\mathcal{M}}(A)| = \bigcup \overline{X} \quad \text{pour } X \in \mathcal{M} \text{ t.q.\ } |X| \subset P \smallsetminus A
\]
\note{« $\overline{X}$ » est sur la page, pour $|X|$ ; de même « $P \smallsetminus A$ » pour $P \smallsetminus |A|$ ; un $\mathcal{M}$ est écrit en surcharge dans « $X \in \mathcal{M}$ »}
d'où l'inclusion dite, sans garantie qu'on soit tjrs égalité.

On aimerait aussi avoir, pour $A$ partie quelconque de $\mathcal{M}$
\[
  (**) \qquad |A| \overset{?}{=} |\operatorname{supp} A| ,
\]
alors qu'on a a priori (tautologiquement)
\[
  |A| \subset |\operatorname{supp} A|
\]
puisque $A \subset \operatorname{supp} A$. Ce \uncertain{dont} on \uncertain{a} \uncertain{besoin} :
\[
  |\operatorname{supp} A| \subset |A|
\]
ce qui serait conséquence de l'implication inverse (hypothétique) de celle du Lemme. La validité de cette implication inverse, pour \struck{\ill{}} $A$, $B$, est d'ailleurs équivalente au cas particulier d'un couple $(\operatorname{supp} A, A)$, i.e.\ à la validité de (**). Cela fait donc quelques perplexités, liées au formalisme \struck{\ill{}} des $P$-supports.
\note{« (**) » : les deux astérisques sont lus sans certitude ; dans la marge de la dernière ligne, un petit triangle surchargé, relié par un trait à un « (2) » cerclé}

Tout se résoud si on fait l'hypothèse suivante

\emph{Hyp} \quad Soit $A \subset \mathcal{M}$, et $x \in P \smallsetminus |A|$. Alors $\exists\, X \in \mathcal{M}$ avec \struck{\ill{}} $X \parallel A$ (i.e.\ $|X| \subset P \smallsetminus |A|$) et $x \in |X|$
\note{« Hyp » est souligné}

\page{78}
\note{son numéro « 78 » en haut de la page}
Mais \uncertain{ça} signifie ni plus ni moins que
\[
  \bigcup_{\substack{X \in \mathcal{M} \\ X \parallel A}} |X| = P \smallsetminus |A|
  \qquad (\text{i.e.\ } X \in \operatorname{cosupp} A)
\]
i.e.\ que l'inclusion $|\complement_{\mathcal{M}}(A)| \subset P \smallsetminus |A|$ est une égalité. (En l'occurrence pas seulement pour $A$ \struck{\ill{}} un ensemble support, mais quelconque…) Comme $\operatorname{supp} A = \operatorname{cosupp}(\operatorname{cosupp} A)$ on en déduira bien, $|\operatorname{supp} A| = P \smallsetminus (P \smallsetminus |A|) = A$, OK.
\marginal{$|\operatorname{cosupp} A| = P - A$}
\note{note écrite en oblique dans la marge gauche, cerclée, avec une flèche vers « i.e.\ » ; « $= A$ » est sur la page, pour $|A|$, comme « $P - A$ » dans la marge. Le premier « $\operatorname{cosupp}$ » de $\operatorname{cosupp}(\operatorname{cosupp} A)$ est écrit en surcharge ; sous « OK. », un mot biffé}

\emph{Lemme 2} \quad Sous les conditions (iii$_{\mathcal{M}}$) et (vi$_{\mathcal{M}}$ a), supposons de plus l'hypothèse précédente satisfaite, i.e.
\[
  \forall A \subset \mathcal{M}, \qquad |\operatorname{cosupp} A| = P \smallsetminus |A| .
\]
Alors pour \uncertain{tout} $A \subset \mathcal{M}$, on a
\[
  |A| = |\operatorname{supp} A| ,
\]
et pour $A, B \subset \mathcal{M}$, on a
\[
  \operatorname{supp} A \subset \operatorname{supp} B \iff |A| \subset |B|
\]
\[
  (\ \iff A \subset \operatorname{supp} B \iff \operatorname{cosupp} B \subset \operatorname{cosupp} A
\]
\note{« Lemme 2 » est souligné. Sur la page, les deux dernières conditions sont écrites en colonne sous « $\operatorname{supp} A \subset \operatorname{supp} B$ », reliées par des signes $=$ verticaux, dans une parenthèse ouverte que ferme la ligne suivante}
et pas seulement « $\Leftarrow$ » comme dans le lemme 1). \uncertain{Ainsi} (\uncertain{cor.}\ \uncertain{lemme}), l'application (injection, plongement ordonné)
\[
  \Sigma_{\mathcal{M}} \longrightarrow \mathfrak{P}(P) \qquad A \longmapsto |A|
\]
commute : l'application $\complement_{\mathcal{M}}$, $\complement_{P}$.

\page{79}
\note{son numéro « 79 » en haut de la page}
Qu'en \uncertain{est}-il des formules pour les Sup et les Inf ? Moyennant passage aux complémentaires, on est réduit à celle sur les Inf, qui est l'intersection.
\[
  (\times) \qquad \struck{\ill{}}\ \left| \bigcap A_i \right| \overset{?}{=} \bigcap_i |A_i|
\]
(les $A_i$ des parties supports).
\marginal{En dehors des cas ensembliste et quasi-ensembliste, ces sup ne sont \uncertain{guère} \uncertain{intéressants}, n'ayant d'intérêt que pour \uncertain{des} parties $A_i$ formées d'éléments \uncertain{deux à deux} compatibles…}
\note{note écrite en oblique dans la marge gauche, en regard du début de la page et de $(\times)$, avec une flèche vers « Sup et les Inf » ; « $(\times)$ » y est répété}

\struck{Pour \uncertain{prendre} le complément des premiers, soit $\mathcal{U}$ \ill{} \ill{} il \uncertain{faut} montrer $\mathcal{U} \cap \bigcap |A_i| = \emptyset$ ?? On pose \uncertain{\ill{}} $\mathcal{U} = |\bigvee A_i|$}
\note{ces lignes sont barrées de grands traits obliques croisés ; la page porte « $\bigvee_i$ » dans la dernière formule biffée}

Je ne sais pas aller plus loin, sauf si je fais l'hypothèse \uncertain{très} commode suivante, beaucoup plus forte que la précédente :

\emph{Hyp'} \quad $\forall\, a \in P$, $\exists\, X \in \mathcal{M}$ avec $|X| = \lbrace a \rbrace$.
\note{« Hyp' » est souligné ; le $P$ de « $a \in P$ » est écrit en surcharge}

Cela implique déjà la surjection (vi$_{\mathcal{M}}$ b), la compatibilité de $\Sigma_{\mathcal{M}} \to \mathfrak{P}(P)$ aux applications $\complement$ (plus gén.\ la formule
\[
  |\operatorname{cosupp} A| = P \smallsetminus |A| \qquad \forall A \subset \mathcal{M} ,
\]
et je dis que ça implique aussi l'égalité \struck{dans} \uncertain{dans} $|X|$ ci-dessus. Soit en effet un pt \uncertain{du} second membre, \struck{\ill{}} soit $X \in \mathcal{M}$ \add{\uncertain{\ill{}}} avec $|X| = \lbrace x \rbrace$, on a \add{$\forall i$} \struck{\ill{}} $|X| \subset |A_i| \iff \operatorname{supp} X \subset \operatorname{supp} A_i$
\note{« surjection » : lecture douteuse ; « $|X|$ ci-dessus » renvoie sans doute à la formule $(\times)$. Sous le dernier $\iff$, deux points lus comme « $A_i$ » ; la page s'arrête sur cette formule}

\page{80}
\note{son numéro « 80 » en haut de la page}
i.e.\ $X \in A_i$, donc $X \in \bigcap A_i$, donc $|X| \subset |\bigcap A_i|$ d'où $x \in |\bigcap A_i|$.

On trouve ainsi

\emph{Proposition} \quad Soit $\mathcal{M}$ un strates, \struck{\ill{}} $P$ un ens., et
\[
  X \longmapsto |X| : \mathcal{M} \longrightarrow \mathfrak{P}(P)
\]
satisfaisant les conditions suivantes :

\emph{Supp}$_{P}^{1}$ \quad \struck{\ill{}} $\forall X, Y \in \mathcal{M}$, $X \parallel Y \iff |X| \cap |Y| = \emptyset$

\emph{Supp}$_{P}^{2}$ \quad $\forall x \in P$, $\exists\, X \in \mathcal{M}$ t.q.\ $|X| = \lbrace x \rbrace$
\note{« Supp$_P^1$ » et « Supp$_P^2$ » sont soulignés ; devant « $\forall X, Y \in \mathcal{M}$ », une formule est raturée de hachures ; dans Supp$_P^2$ les lettres $P$ et $\mathcal{M}$ sont écrites en surcharge}

Posons, pour tout $A \subset \mathcal{M}$,
\[
  |A| \overset{\mathrm{déf}}{=} \bigcup_{X \in A} |X| .
\]
\emph{Alors} on a ce qui suit

(i) $\forall A \subset \mathcal{M}$, on a
\[
  |A| = \operatorname{supp} A .
\]
\note{la page écrit « $|A| = \operatorname{supp} A$ », sans barres au second membre ; le lemme 2 (page 78) énonce $|A| = |\operatorname{supp} A|$}
\struck{(ii) Pour} $|\operatorname{cosupp} A| = P \smallsetminus |A|$.
\note{cette formule est écrite sous la précédente, après un « (ii) Pour » biffé}

(ii) Pour $A, B \subset \mathcal{M}$, on a

\quad 1°) $\operatorname{supp} A \subset \operatorname{supp} B$ \add{i.e.\ $A \subset \operatorname{supp} B$} $\iff |A| \subset |B|$

\quad 2°) $A \parallel B \iff |A| \cap |B| = \emptyset$

(iii) L'application
\[
  \Sigma_{\mathcal{M}} \rightleftarrows \mathfrak{P}(P) \qquad A \longmapsto |A|
\]
\struck{est bijective} \uncertain{est} \add{un} iso.\ d'ens.\ ordonnés commutant aux op.\ $\complement$.
\marginal{Ceci est valable sans Supp$^{2}$ et \uncertain{avec} \uncertain{moins} \uncertain{même}}
\note{note dans la marge gauche, avec une flèche vers 2°). Le numéro de (iii) est écrit en surcharge ; « Alors » est souligné. La page, et le lot, s'arrêtent sur cet énoncé}

\end{document}
